\section*{Chapter \ref{ch:b3:exchange-rate-regimes} --- Exchange-Rate Regimes and Currency Unions}

\begin{solution}{exo:b3:exchange-rate-regimes:1}
The \omterm{def:b3:exchange-rate-regimes:gold-standard}{gold standard} (from the 1870s to 1914; 82~\% pegged in 1880 and 1913); its
collapse in the 1930s (39~\% in 1932); the Bretton Woods system (1950s and 1960s,
100~\% in 1970); and the floats, European \omterm{def:b3:exchange-rate-regimes:regimes}{pegs} and the euro since the 1970s (about
50 to 60~\%). The share was highest under Bretton Woods. % figdata: pegs_points
\end{solution}

\begin{solution}{exo:b3:exchange-rate-regimes:2}
(a) Fixed rate and free capital movements, without its own monetary policy. (b)
Free capital and its own policy, without a fixed rate. (c) Fixed rate and its own
policy, without free capital. (d) Fixed rate and free capital, without its own
policy.
\end{solution}

\begin{solution}{exo:b3:exchange-rate-regimes:3}
$\tilde S = 40/50 + 100 \times 2/50^2 = 0.80 + 0.08 = 0.88$. Below the \omterm{def:b3:exchange-rate-regimes:regimes}{peg} of 1.0: an attack would
leave speculators holding foreign currency worth less than they paid, so the \omterm{def:b3:exchange-rate-regimes:regimes}{peg}
is not under attack yet.
\end{solution}

\begin{solution}{exo:b3:exchange-rate-regimes:4}
A \omterm{def:b3:exchange-rate-regimes:regimes}{peg} that markets believe and capital that moves freely, with home interest
rates that move independently of the base country's for years --- a slope near
zero in the trilemma regression --- without reserve losses or an expected
depreciation to account for the difference.
\end{solution}

\begin{solution}{exo:b3:exchange-rate-regimes:5}
The tested solution is \code{ex05_trilemma.py}, in the chapter's \code{exercises/}
folder. Slopes, \omterm{def:b3:exchange-rate-regimes:regimes}{pegs} then floats: 1870--1913, 0.41 and 0.52; 1950--1972, 0.28 and 0.53;
1973--1998, 0.45 and 0.37; 1999--2020, 0.91 and 0.45. % output: ex05_trilemma
Since 1999 the contrast is the one the trilemma predicts. Under Bretton Woods the
\omterm{def:b3:exchange-rate-regimes:regimes}{pegs} followed the base least, which makes sense only with the third corner:
capital controls let pegged countries keep their own rates. Before 1914 the floats
are few (150 observations) and their estimate imprecise; from 1973 to 1998 the \omterm{def:b3:exchange-rate-regimes:regimes}{pegs}
included realigned European \omterm{def:b3:exchange-rate-regimes:regimes}{pegs} and capital controls lasted into the period.
\end{solution}

\begin{solution}{exo:b3:exchange-rate-regimes:6}
$T = R_0/\mu - c/k = 40/4 - 2 = 8$ years; reserves lost $c\mu/k = 200 \times 4/100 = 8$. With $c = 400$:
$T = 10 - 4 = 6$, and 16 are lost. A larger $c$ means money demand falls more when the
interest rate jumps at the attack, so the attack takes more reserves and must come
while more remain.
\end{solution}

\begin{solution}{exo:b3:exchange-rate-regimes:7}
With 14: even both together (12) cannot exhaust reserves; selling always costs 0.5,
so holding is dominant: only (hold, hold). With 5: one seller alone (6) breaks the
\omterm{def:b3:exchange-rate-regimes:regimes}{peg} and gains 1; selling is dominant: only (sell, sell). More reserves (or credit
lines), a credible willingness to defend, sounder banks that make high interest
rates less costly, and capital controls all move a country toward the first case.
\end{solution}

\begin{solution}{exo:b3:exchange-rate-regimes:8}
The tested solution is \code{ex08_greece_germany.py}, in the chapter's
\code{exercises/} folder. The gap ranged from 0.7 points (2007) to 2.6 (2002), with a
mean of 1.60 points; over the ten years Greek prices rose 16.9~\% relative to German
ones. % output: ex08_greece_germany
\end{solution}

\begin{solution}{exo:b3:exchange-rate-regimes:9}
Endogenous: correlations of members' cycles that rise after they integrate or join
a union, more than those of comparable non-members; trade among members that
grows after the union. Not endogenous: correlations unchanged or falling after
integration, as when trade leads countries to specialise in different industries
and so to face different shocks.
\end{solution}

\begin{solution}{exo:b3:exchange-rate-regimes:10}
After the attack the rate follows the shadow rate. If the attack came when $\tilde S <
\bar S$, the rate would fall at once to $\tilde S$: speculators who bought kests at $\bar S$ would lose.
If it came when $\tilde S > \bar S$, the rate would jump up, a sure gain at an infinite rate of
return, and speculators would have attacked a moment earlier to take it. The only
date consistent with no expected loss and no unexploited gain is $\tilde S = \bar S$: the
attack comes then, and the rate moves continuously.
\end{solution}

\begin{solution}{exo:b3:exchange-rate-regimes:11}
Relative wages must fall by 15~\%: $(1/1.02)^n = 0.85$, $n = \ln(1/0.85)/\ln 1.02 = 0.1625/0.0198 =
8.2$ years with a freeze; with wages falling by 1~\%, $(0.99/1.02)^n = 0.85$, $n = 0.1625/0.0299
= 5.4$ years.
\end{solution}

\begin{solution}{exo:b3:exchange-rate-regimes:12}
The base falls to 92, and bank reserves to 12. To keep reserves at 10~\% of deposits,
banks shrink deposits to 120: a fall of 80, of which 8 is the outflow and 72 the
loans called in. A central bank would have lent to the banks against collateral,
keeping their reserves and the base from falling --- the lender of last resort a
board cannot be.
\end{solution}

\begin{solution}{pb:b3:exchange-rate-regimes:peg}
\textbf{1.} $M = k\bar S = 100 \times 1.20 = 120$ billion; $R_0 = 120 - 90 = 30$ billion dalens. % figdata: attack_stats, attack

\textbf{2.} Money demand is fixed at 120 under the \omterm{def:b3:exchange-rate-regimes:regimes}{peg}, so the new credit is money
nobody wants to hold: it is exchanged for kests at the central bank. $R_t = 30 - 3t$.

\textbf{3.} In 10 years.

\textbf{4.} $\tilde S_0 = 90/100 + 200 \times 3/100^2 = 0.90 + 0.06 = 0.96$, below the \omterm{def:b3:exchange-rate-regimes:regimes}{peg}: no pressure.

\textbf{5.} $\tilde S_t = (90 + 3t)/100 + 0.06 = 0.96 + 0.03t$.

\textbf{6.} $30/1.20 = 25$ billion kests.

\textbf{7.} $0.96 + 0.03T = 1.20$: $T = 8$ years ($= 30/3 - 200/100$).

\textbf{8.} $c\mu/k = 6$ billion dalens.

\textbf{9.} $120 - 6 = 114$ billion: $k\bar S - c\dot S = 120 - 200 \times 0.03$.

\textbf{10.} $\dot S = \mu/k = 0.03$ dalens per kest a year, $0.03/1.20 = 2.5~\%$ a year at first.

\textbf{11.} By the expected depreciation, 2.5 points, under uncovered parity.

\textbf{12.} The attack comes exactly when the floating rate would equal the \omterm{def:b3:exchange-rate-regimes:regimes}{peg};
any other date would offer a loss or an unexploited gain.

\textbf{13.} $T = 30/6 - 2 = 3$ years.

\textbf{14.} $R_4 = 18$ billion; $\tilde S_4 = 1.08 < 1.20$: no.

\textbf{15.} Each fund's gain if the \omterm{def:b3:exchange-rate-regimes:regimes}{peg} falls: $0.20 \times 10 - 0.4 = 1.6$. Together they sell
20 > 18 and the \omterm{def:b3:exchange-rate-regimes:regimes}{peg} falls; one alone sells 10 < 18 and it holds. (Sell, sell): 1.6,
1.6; (sell, hold): $-0.4$, 0; (hold, sell): 0, $-0.4$; (hold, hold): 0, 0.

\textbf{16.} Both sell, and both hold: two equilibria.

\textbf{17.} The government gives up after $18 - 12 = 6$ billion of sales, less than one
fund's 10: a single fund breaks the \omterm{def:b3:exchange-rate-regimes:regimes}{peg} and gains 1.6. Selling is dominant; the
only equilibrium is an attack.

\textbf{18.} In Question 16 the \omterm{def:b3:exchange-rate-regimes:regimes}{peg} is sound (the shadow rate is 1.08), yet it falls
if each fund expects the other to sell: the expectation creates the devaluation
that justifies it. The threshold of Question 17 removes the good equilibrium: what
the government will bear, not its reserves, decides.

\textbf{19.} Reserves above $12 + 20 = 32$ billion (a credit line from abroad), so that even a
joint attack fails; a credible commitment to defend at any cost; slower credit
growth, which also delays the first-generation attack; or capital controls.

\textbf{20.} The full base, 120 billion dalens: $120/1.20 = 100$ billion kests.

\textbf{21.} $30/120 = 25~\%$; the shortfall is 90 billion dalens, 75 billion kests.

\textbf{22.} The board cannot issue base money against government debt: the deficit
must be financed by taxes or by borrowing from markets.

\textbf{23.} The base falls to 108 billion, 10~\% less. Money demand at the \omterm{def:b3:exchange-rate-regimes:regimes}{peg} is still
120: interest rates rise until spending and the demand for money fall to match
--- a monetary contraction no one decided.

\textbf{24.} The attack comes in year 8, when 6 billion dalens of reserves remain; a
\omterm{def:b3:exchange-rate-regimes:regimes}{currency board} would need 120 billion dalens of reserves (100 billion kests), four
times what Brenland holds.
\end{solution}
